\chapter{镜像下降的几何、策略改进与参数求解}
\label{ch:mirror-foundations}

镜像下降先在概率分布空间中定义目标；策略镜像下降（policy mirror descent，PMD）把目标接到当前策略的价值函数；参数化优化器再把目标映射到神经网络能够表达的策略。自然梯度（natural gradient，NG）描述第三层中的局部距离几何。本章分别建立这些对象，使后续能够明确说明一个 outer step 完成了什么，以及有限 inner steps 留下了什么误差。

\section{Bregman 散度与镜像子问题}

\subsection{势函数和三点恒等式}

设 $\mathcal C$ 是凸集合，$h$ 在包含其相对内部的区域可微、严格凸。由 $h$ 生成的 Bregman 散度为
\begin{equation}
 D_h(x,y)=h(x)-h(y)-\langle\nabla h(y),x-y\rangle.
 \label{eq:foundation-bregman}
\end{equation}
凸性给出 $D_h(x,y)\geq0$；严格凸性给出等号当且仅当 $x=y$。它通常不对称，也不满足三角不等式。镜像下降只需要用它度量候选点偏离旧迭代点的代价。

\begin{lemma}[Bregman 三点恒等式]
\label{lem:foundation-three-point}
在势函数可微区域，对任意 $x,y,z$ 有
\begin{equation}
 \langle\nabla h(x)-\nabla h(y),z-x\rangle
 =D_h(z,y)-D_h(z,x)-D_h(x,y).
 \label{eq:foundation-three-point}
\end{equation}
\end{lemma}
\begin{proof}
分别展开右边：
\begin{align*}
 D_h(z,y)&=h(z)-h(y)-\langle\nabla h(y),z-y\rangle,\\
 D_h(z,x)&=h(z)-h(x)-\langle\nabla h(x),z-x\rangle,\\
 D_h(x,y)&=h(x)-h(y)-\langle\nabla h(y),x-y\rangle.
\end{align*}
相减后函数值全部抵消，$\nabla h(y)$ 的两项合并为
\[
 -\langle\nabla h(y),z-y\rangle+\langle\nabla h(y),x-y\rangle
 =-\langle\nabla h(y),z-x\rangle.
\]
加上剩余的 $\langle\nabla h(x),z-x\rangle$，得到结论。
\end{proof}

最小化方向的镜像步定义为
\begin{equation}
 x^+=\argmin_{u\in\mathcal C}
 \{\eta\langle g,u\rangle+D_h(u,x)\},\qquad\eta>0.
 \label{eq:foundation-md-min}
\end{equation}
$g$ 是当前目标梯度或其估计。$\eta$ 控制目标分布相对参考点的移动，而 inner 参数优化器学习率控制求解过程，二者有不同语义。

\begin{proposition}[镜像步的三点不等式]
\label{prop:foundation-md-three-point}
若解存在且 $h$ 在 $x^+$ 可微，则任意 $u\in\mathcal C$ 满足
\begin{equation}
 \eta\langle g,x^+-u\rangle
 \leq D_h(u,x)-D_h(u,x^+)-D_h(x^+,x).
 \label{eq:foundation-md-three-point}
\end{equation}
\end{proposition}
\begin{proof}
凸子问题的一阶最优性条件为
\[
 \langle\eta g+\nabla h(x^+)-\nabla h(x),u-x^+\rangle\geq0.
\]
移项得
\[
 \eta\langle g,x^+-u\rangle
 \leq\langle\nabla h(x^+)-\nabla h(x),u-x^+\rangle.
\]
在引理 \ref{lem:foundation-three-point} 中取 $(x,y,z)=(x^+,x,u)$ 即得结论。
\end{proof}

这条估计使相邻轮次的散度能够望远镜相消。其证明没有使用 MDP 结构，因此它先是通用凸优化工具，随后才被接到策略改进上。

\subsection{凸目标的平均误差界}

\begin{theorem}[凸镜像下降的误差界]
\label{thm:foundation-convex-md}
设 $\phi$ 为凸函数，$g_k\in\partial\phi(x_k)$，$\lVert g_k\rVert_*\leq G$，并且
\[
 D_h(u,v)\geq\frac\sigma2\lVert u-v\rVert^2,\qquad\sigma>0.
\]
用固定 $\eta>0$ 执行镜像步。对最优解 $x^*$ 和 $K\geq1$，有
\begin{equation}
 \phi\left(\frac1K\sum_{k=0}^{K-1}x_k\right)-\phi(x^*)
 \leq\frac{D_h(x^*,x_0)}{K\eta}+\frac{\eta G^2}{2\sigma}.
 \label{eq:foundation-convex-md-rate}
\end{equation}
\end{theorem}
\begin{proof}
由子梯度不等式和三点估计，
\begin{align*}
 \eta[\phi(x_k)-\phi(x^*)]
 &\leq\eta\langle g_k,x_k-x_{k+1}\rangle
      +\eta\langle g_k,x_{k+1}-x^*\rangle\\
 &\leq D_h(x^*,x_k)-D_h(x^*,x_{k+1})
      -D_h(x_{k+1},x_k)+\eta G\lVert x_k-x_{k+1}\rVert.
\end{align*}
令 $r=\lVert x_k-x_{k+1}\rVert$，则
\[
 -\frac\sigma2r^2+\eta Gr
 =-\frac\sigma2\left(r-\frac{\eta G}\sigma\right)^2
   +\frac{\eta^2G^2}{2\sigma}
 \leq\frac{\eta^2G^2}{2\sigma}.
\]
对 $k$ 求和，散度项相消，得到
\[
 \eta\sum_{k=0}^{K-1}[\phi(x_k)-\phi(x^*)]
 \leq D_h(x^*,x_0)-D_h(x^*,x_K)+\frac{K\eta^2G^2}{2\sigma}.
\]
删除非正项，除以 $K\eta$，再用凸性处理迭代平均即可。
\end{proof}

固定步长意味着 $\eta$ 不随 $k$ 改变，梯度、参考点和每轮数据依然更新。策略参数目标一般非凸；PMD 的全局分析利用 Bellman 结构，不能直接把本定理中的 $x$ 换成 Transformer 参数。

\section{负熵镜像与 Gibbs 变分恒等式}

在概率单纯形上取负熵 $h(p)=\sum_ap(a)\log p(a)$。对严格正的 $q$，$\nabla h(q)_a=\log q(a)+1$。利用概率总和为 $1$，有
\[
 D_h(p,q)=\sum_ap(a)\log\frac{p(a)}{q(a)}=\KL(p\Vert q).
\]
若 $q$ 有零概率，候选必须限制在 $\supp(q)$ 上；否则散度为无穷大。

固定状态 $s$，暂时省略 $s$，设 anchor 为 $\mu$，冻结信号为 $f$。镜像步长 $\eta$ 与温度 $\tau$ 的关系为 $\eta=1/\tau$：
\begin{equation}
 J_\mu(p)=\E_p[f]-\frac1\eta\KL(p\Vert\mu),\qquad
 p^+=\argmax_{p\ll\mu}J_\mu(p).
 \label{eq:foundation-gibbs-objective}
\end{equation}

\begin{theorem}[Gibbs 解和精确目标差]
\label{thm:foundation-gibbs}
设动作集合有限、$\eta>0$、$f$ 在 $\supp(\mu)$ 上有限。定义
\[
 Z=\sum_a\mu(a)e^{\eta f(a)},\qquad
 p^*(a)=\frac{\mu(a)e^{\eta f(a)}}Z.
\]
则 $p^*$ 为唯一最大化解。任意 $p\ll\mu$ 满足
\begin{equation}
 J_\mu(p)=\frac{\log Z}\eta-\frac1\eta\KL(p\Vert p^*),
 \qquad J_\mu(p^*)-J_\mu(p)=\frac1\eta\KL(p\Vert p^*).
 \label{eq:foundation-gibbs-gap}
\end{equation}
\end{theorem}
\begin{proof}
有限动作和有限信号保证 $0<Z<\infty$。$p^*$ 非负、总和为 $1$，并与 $\mu$ 有相同 support。由定义逐动作有
\[
 \log\frac{p(a)}{p^*(a)}
 =\log\frac{p(a)}{\mu(a)}-\eta f(a)+\log Z.
\]
在 $p$ 下求期望，得到
\[
 \KL(p\Vert p^*)=\KL(p\Vert\mu)-\eta\E_p[f]+\log Z.
\]
整理得式 \eqref{eq:foundation-gibbs-gap}。KL 非负且仅在 $p=p^*$ 时为零，故唯一最优性和目标差同时成立。
\end{proof}

拉格朗日方法给出的内点条件为
\[
 f(a)-\frac1\eta\left(\log\frac{p(a)}{\mu(a)}+1\right)+\lambda=0.
\]
消去乘子同样得到 Gibbs 形式。变分恒等式额外给出全局最优性及目标差与 KL 的精确对应，因而适合分析有限 inner 求解误差。

将 $f$ 减去动作无关的常数 $c(s)$ 不改变 target，因为分子分母同时乘 $e^{-\eta c(s)}$。因此 $Q^\mu$ 和 $A^\mu=Q^\mu-V^\mu$ 给出同一状态级 target。均值为零却不意味着配分函数为 $1$：若 $\E_\mu f=0$，Jensen 不等式给出 $Z\geq1$，非恒定信号时严格大于 $1$。

\section{PMD 与 Bellman 策略改进}

\subsection{精确无正则镜像步}

考虑有限状态、有限动作、折扣 $\gamma\in[0,1)$ 的 MDP，采用奖励最大化约定。动作价值 $Q^\pi(s,a)$ 表示先执行 $a$、随后遵循 $\pi$ 的期望折扣回报，且
\[
 V^\pi(s)=\langle Q^\pi(s,\cdot),\pi(\cdot\mid s)\rangle,
 \qquad A^\pi(s,a)=Q^\pi(s,a)-V^\pi(s).
\]
一轮精确 KL-PMD 为
\begin{equation}
 \pi_{k+1}(\cdot\mid s)=\argmax_p
 \left\{\langle Q^{\pi_k}(s,\cdot),p\rangle
       -\frac1{\eta_k}\KL(p\Vert\pi_k(\cdot\mid s))\right\}.
 \label{eq:foundation-pmd}
\end{equation}
表格、严格正旧策略和精确价值函数条件下，每个状态由 Gibbs 公式独立更新。算法本身没有必须运行到收敛的数值 inner loop；inner loop 来自近似表示或近似价值评估。

\begin{proposition}[精确 PMD 的一轮策略改进]
\label{prop:foundation-pmd-improvement}
精确执行式 \eqref{eq:foundation-pmd} 时，每个状态满足
\[
 \langle A^{\pi_k},\pi_{k+1}\rangle
 \geq\frac1{\eta_k}\KL(\pi_{k+1}\Vert\pi_k)\geq0,
\]
并且逐状态有 $V^{\pi_{k+1}}\geq V^{\pi_k}$。
\end{proposition}
\begin{proof}
将 $\pi_k$ 本身作为可行候选，最优性给出
\[
 \langle Q^{\pi_k},\pi_{k+1}-\pi_k\rangle
 \geq\frac1{\eta_k}\KL(\pi_{k+1}\Vert\pi_k).
\]
左边恰等于 $\langle A^{\pi_k},\pi_{k+1}\rangle$。将这个非负状态向量记为 $b$，新策略的转移矩阵记为 $P_{k+1}$。两条 Bellman 方程相减得
\[
 V^{\pi_{k+1}}-V^{\pi_k}=b+\gamma P_{k+1}(V^{\pi_{k+1}}-V^{\pi_k}).
\]
从而
\[
 V^{\pi_{k+1}}-V^{\pi_k}
 =(I-\gamma P_{k+1})^{-1}b
 =\sum_{j=0}^\infty\gamma^jP_{k+1}^jb\geq0.
\]
级数在 $\gamma<1$ 时收敛，每项非负，结论成立。
\end{proof}

本证明使用当前策略的完整动作价值。将信号换成即时 token reward 或教师 log-ratio 后，Gibbs 定理依然有效，Bellman 改进部分则需要另行建立信号与价值的关系。

\subsection{一般正则项的三点估计}

这一小节采用代价最小化约定，设 $q_k$ 为当前动作代价，$\Omega_s$ 为凸策略正则项。子问题为
\[
 p^+=\argmin_p\{\eta[\langle q_k,p\rangle+\Omega_s(p)]+D_h(p,p_k)\}.
\]
若 $\Omega_s$ 相对于 $h$ 为 $\mu$ 强凸，即对 $v\in\partial\Omega_s(p^+)$ 有
\[
 \Omega_s(p)\geq\Omega_s(p^+)+\langle v,p-p^+\rangle+\mu D_h(p,p^+),
\]
则有
\begin{equation}
 \eta[\langle q_k,p^+-p\rangle+\Omega_s(p^+)-\Omega_s(p)]
 \leq D_h(p,p_k)-(1+\eta\mu)D_h(p,p^+)-D_h(p^+,p_k).
 \label{eq:foundation-regularized-pmd}
\end{equation}
证明时先用一阶条件和三点恒等式控制 $\eta\langle q_k+v,p^+-p\rangle$。再由强凸性得到
\[
 \Omega_s(p^+)-\Omega_s(p)\leq\langle v,p^+-p\rangle-\mu D_h(p,p^+).
\]
将后一式乘以 $\eta$，与前一式合并，即得到额外的 $-\eta\mu D_h(p,p^+)$。

Lan 的 PMD 分析把状态级估计与正则化 Bellman 方程、占用测度结合，在规定步长与正则条件下得到全局速率。原文固定步长线性收敛定理有强凸条件；一般凸情形和 approximate PMD 使用各自的构造与参数。正则化 PMD 文献还单独处理不精确价值评估和不精确策略更新。式 \eqref{eq:foundation-regularized-pmd} 是本书直接证明的代数工具，不能据此宣称任意神经 OPD 满足论文的全局定理。\cite{lan23,zhan2022}

\subsection{固定步长与每轮刷新}

精确 PMD 每轮先给定 $\pi_k$，评估 $Q^{\pi_k}$；固定这两个输入，逐状态执行 prox map；得到 $\pi_{k+1}$ 后重新评估新策略的动作价值。常数 $\eta_k=\eta$ 固定系数，旧策略和动作价值仍随 outer 更新。随机估计每轮的对象仍是 $Q^{\pi_k}$。永久复用 $\pi_0$ 的价值估计，就变成固定 surrogate 上的优化，改变了策略迭代的对象。

\section{MDPO 的 on-policy 与 off-policy pipeline}

\subsection{On-policy：固定参考策略的多步求解}

Tomar 等人的 MDPO 用参数化策略近似求解旧访问测度下的镜像子问题。其 on-policy 目标为\cite{mdpo20}
\begin{equation}
 \Psi_k(\theta)=\E_{s\sim d_{\pi_k}}
 \left[\E_{a\sim\pi_\theta}A^{\pi_k}(s,a)
 -\frac1{t_k}\KL(\pi_\theta(\cdot\mid s)\Vert\pi_k(\cdot\mid s))\right].
 \label{eq:foundation-mdpo}
\end{equation}
这里保留原文步长符号 $t_k$。它是 prox 系数，与参数优化器学习率不同。用旧策略数据估计第一项时，写成
\[
 \E_{s\sim d_{\pi_k},a\sim\pi_k}
 \left[\frac{\pi_\theta(a\mid s)}{\pi_k(a\mid s)}\widehat A_k(s,a)\right].
\]
这个动作比率校正固定状态下的动作期望；状态测度仍为 $d_{\pi_k}$，与原来的 surrogate 一致。

原文 Algorithm 1 的更新次序是：
\begin{enumerate}
 \item 用当前 $\pi_k$ 生成新的 rollout，计算 return 与优势，并更新价值网络。
 \item 保存 old-policy 参数 $\theta_k$，将 candidate 初始化为 $\theta_{k,0}=\theta_k$。
 \item 固定 rollout、优势与 old-policy copy，对 $\Psi_k$ 执行 $m$ 次 SGD。
 \item 令 $\theta_{k+1}=\theta_{k,m}$，刷新 old copy，再进行下一轮 rollout。
\end{enumerate}
原文连续控制实验比较了不同 $m$，on-policy 采用 $m=10$。这是实验中的近似求解预算，不是 outer 次数，也不是 LLM 的通用停止标准。

\begin{proposition}[Anchor 处 KL 的梯度为零]
\label{prop:foundation-kl-zero}
对 support 固定、可微且允许交换求导与求和的概率族，定义 $K_k(\theta)=\KL(\pi_\theta\Vert\pi_{\theta_k})$，则
\[
 \nabla K_k(\theta_k)=0.
\]
从 $\theta_k$ 开始仅做一次普通显式梯度上升，KL 惩罚不影响该次一阶方向。
\end{proposition}
\begin{proof}
逐项求导得
\[
 \nabla K_k(\theta)=\sum_a\nabla\pi_\theta(a)
 \left(\log\frac{\pi_\theta(a)}{\pi_{\theta_k}(a)}+1\right).
\]
在 $\theta_k$ 处括号为 $1$，故
\[
 \nabla K_k(\theta_k)=\nabla\sum_a\pi_\theta(a)\big|_{\theta_k}=0.
\]
连续密度把求和换成积分即可。显式一步更新中的惩罚梯度因此为零。
\end{proof}

这一结论针对一阶更新：二阶信赖域法仍可在 anchor 处使用 KL 曲率。普通 SGD 后续步骤因为离开了 anchor，KL 梯度开始生效。若每步都同步 anchor，梯度归零会反复出现，与固定子问题的多步求解产生不同动力学。

\subsection{Off-policy：replay 状态与重新采样的动作}

MDPO 的 off-policy 推导从 replay buffer 取状态，近似较均匀的状态采样；原文明确说明这不是真正均匀分布。其 actor 子问题为
\[
 \E_{s\sim\mathcal D}
 \left[\E_{a\sim\pi_\theta}Q_{\psi,k}(s,a)
       -\frac1{t_k}\KL(\pi_\theta\Vert\pi_k)(s)\right].
\]
连续动作通过 $a=\widetilde a_\theta(\epsilon,s)$ 重参数化，等价 actor 损失为
\[
 \E_{s\sim\mathcal D,\epsilon}
 [\log\pi_\theta(\widetilde a_\theta\mid s)
 -\log\pi_k(\widetilde a_\theta\mid s)-t_kQ_{\psi,k}(s,\widetilde a_\theta)].
\]
actor 动作由 candidate 重新采样，因此这里没有使用 replay 存下来的行为动作再乘 IS 比率。replay transition 用于 TD critic 学习，critic 误差与 replay 状态覆盖仍需单独控制。

原文 Algorithm 2 描述固定 old policy 的多步 SGD。实际 off-policy 实验为效率每次只做一个 actor 梯度更新，将 old copy 保留 $m$ 次更新后再同步，从而模拟多步求解。论文报告其 off-policy 实验采用 $m=1000$；此处 $m$ 是 old copy 的更新周期，不能与 on-policy 的 $m=10$ 直接比较。

MDPO 和 soft actor-critic（SAC）的差别来自 anchor：前者的 Boltzmann target 为 $\pi_k\exp(t_kQ)$，后者的熵正则目标相当于相对均匀基准测度的 Boltzmann 分布。它们都做 KL 投影，一个保留历史策略 reference，一个使用熵正则基准。相同的投影名称没有消除算法状态的差别。

\section{Fisher 几何与自然梯度}

\subsection{KL 展开和可辨识子空间}

固定状态测度 $d$，设策略族严格正、三次连续可微。定义 score 与 Fisher：
\[
 u_\theta(s,a)=\nabla_\theta\log\pi_\theta(a\mid s),\qquad
 F_d(\theta)=\E_{d,\pi_\theta}[u_\theta u_\theta^\top].
\]
概率归一化求导给出 $\E_{\pi_\theta}u_\theta=0$，再求导得
\[
 \E_{\pi_\theta}\nabla^2\log\pi_\theta
 =-\E_{\pi_\theta}[u_\theta u_\theta^\top].
\]
展开新策略 log-prob，旧到新 KL 为
\begin{align*}
 \E_d\KL(\pi_\theta\Vert\pi_{\theta+\delta})
 &=-\E_{d,\pi_\theta}
 [u_\theta^\top\delta+\tfrac12\delta^\top\nabla^2\log\pi_\theta\,\delta]
 +O(\lVert\delta\rVert^3)\\
 &=\tfrac12\delta^\top F_d(\theta)\delta+O(\lVert\delta\rVert^3).
\end{align*}
新到旧 KL 有相同二阶项，但高阶项一般不同。若还改变状态测度，就需要分析额外的测度变化。

Fisher 的核表示当前覆盖测度下没有一阶概率变化的参数方向。在去掉这些方向的商空间求解，或使用 广义逆 $F^\dagger$。必须要求 $g\in\operatorname{range}(F)$；否则梯度有非零零空间分量，仅用 Fisher 半径不能限制该分量，线性目标会无界。

\begin{proposition}[Fisher 信赖域的解]
\label{prop:foundation-natural-trust}
设 $F$ 半正定，$0\neq g\in\operatorname{range}(F)$。在商空间最大化 $g^\top\delta$ 且约束 $\delta^\top F\delta/2\leq\varepsilon$，解为
\[
 \delta^*=\sqrt{\frac{2\varepsilon}{g^\top F^\dagger g}}F^\dagger g.
\]
惩罚问题 $\max_\delta\{g^\top\delta-\delta^\top F\delta/(2\eta)\}$ 的最小范数解为 $\eta F^\dagger g$。
\end{proposition}
\begin{proof}
在正特征值子空间令 $w=F^{1/2}\delta$、$v=F^{\dagger/2}g$。目标等于 $v^\top w$，约束变成 $\lVert w\rVert^2\leq2\varepsilon$。Cauchy--Schwarz 的上界在 $w=\sqrt{2\varepsilon}v/\lVert v\rVert$ 取得，逆变换得第一式。惩罚问题的一阶条件为 $g-F\delta/\eta=0$，给出第二式。最小范数选择把零空间分量设为零。
\end{proof}

这连接 Amari 的信息几何、Kakade 的自然策略梯度和 TRPO 的局部信赖域更新：步长按分布变化度量。阻尼、采样 Fisher 和线性求解残差会改变实际距离，需要检查真实 KL。\cite{amari1998,kakade01,trpo15}

\subsection{Softmax gauge 下的有限步精确恒等式}

令独立 logits $z\in\R^{|\mathcal A|}$，$p=\softmax(z)$。Jacobian 与 logits Fisher 都为
\[
 C(p)=\diag(p)-pp^\top.
\]
因 $\softmax(z+c\1)=\softmax(z)$，有 $C(p)\1=0$。严格正 $p$ 下其核恰为 $\operatorname{span}\{\1\}$。这称为 logits 的 gauge 自由度。

冻结信号向量 $f$，设 $J(z)=p^\top f$，则 $\nabla_zJ=C(p)f$。自然梯度为
\[
 C(p)^\dagger\nabla_zJ=C(p)^\dagger C(p)f
 =f-\frac{\1^\top f}{|\mathcal A|}\1.
\]
显式自然梯度 logits 步因此给出
\[
 p^+=\softmax(z+\eta f),\qquad
 p^+(a)=\frac{p(a)e^{\eta f(a)}}{\sum_bp(b)e^{\eta f(b)}}.
\]
这是有限 $\eta$ 的精确 Gibbs mirror map，常数 gauge 自动消失。因此“一切自然梯度都只与小步镜像下降等价”的表述过强；独立 tabular softmax 对冻结线性信号存在这个有限步恒等式。

该恒等式针对 anchor 处冻结线性目标的一次更新，不意味着固定 KL 正则子问题上的任意自然梯度算法一步收敛。共享参数网络的 logits 无法逐状态独立移动，参数 Fisher 只能把方向投影到模型切空间。

\subsection{一般参数化的局部近端模型}

设 logits Jacobian 为 $J_k(s)$，$\Delta=\theta-\theta_k$，固定测度 $d_k$ 下
\[
 F_k=\E_{d_k}[J_k(s)^\top C(\pi_k(\cdot\mid s))J_k(s)],\qquad
 b_k=\nabla_\theta\E_{d_k,\pi_\theta}f_k\big|_{\theta_k}.
\]
信号目标保留一阶、KL 保留二阶，得到局部模型
\[
 \max_\Delta\{b_k^\top\Delta-\Delta^\top F_k\Delta/(2\eta)\}.
\]
其解 $\eta F_k^\dagger b_k$ 是一般参数族中 mirror 更新的局部表示。真实信号目标也有 Hessian：在 $\Delta=O(\eta)$ 且导数有界时，其影响为更高阶；固定较大 $\eta$ 时不能删除该项后把二次模型当作全局目标。

分布 target 是否可表示、inner 优化器如何求解、局部曲率是否稳定，是三个独立条件。普通 SGD 和 Adam 不直接计算 Fisher 逆；MDPO 多步优化的 endpoint 能逼近镜像子问题，并不说明每个 optimizer step 都是自然梯度。后续分析将二次模型的谱收缩与真实概率残差分开。

\section{文献谱系与 OPD 的接口}

Beck--Teboulle 提供 Bregman prox 和三点估计；信息几何与自然策略梯度提供 KL 的局部 Fisher 度量；TRPO 用半径约束实现局部信赖域；MDPO 用固定旧策略的多步优化逼近镜像子问题；PMD 全局理论利用 Bellman 结构。它们按数学对象连接，不能简化成同一个优化器名称。\cite{beck2003,amari1998,kakade01,trpo15,mdpo20,lan23,zhan2022}

一般参数化 PMD 研究可表示性条件，homotopic PMD 调整正则强度，lookahead、block PMD 与 learned mirror maps 改变更新位置或势函数，随机镜像下降及 Bregman Gradient Policy Optimization 研究随机估计和非凸驻点。Bregman Gradient Policy Optimization 的作者为 Huang、Gao 与 Huang；其参数化驻点保证与 tabular PMD 的最优策略保证属于不同对象。GPOMDP 属于策略梯度估计器谱系，也不能直接据其名称归为自然梯度算法。

REPS、Maximum a Posteriori Policy Optimisation（MPO）和 Guided Policy Search 的共同接口是先构造指数倾斜分布，再投影到参数策略。这里的 MPO 指策略优化方法，偏好混合后处理中的同名缩写属于另一工作。它们启发把镜像温度和学生 Fisher 半径分开，分别控制 target 偏离 anchor 与学生追赶 target 的实际步长。\cite{mpo18}

本章提供三个可直接使用的接口：Gibbs 恒等式确定目标及精确目标差；三点不等式确定固定 anchor 的近端几何；Fisher 公式确定参数空间的局部求解方向。下一章把教师 log-ratio 代入这些接口，建立 OPD 特有的精确收缩。环境回报、访问覆盖与 off-policy 方差还需要各自的信号和测度条件。
